\answer{ex:09:1}{(а)~$R=4$: $P_\infty=0.2$, ошибка $0.45$, скорость $0.05\,\text{с}^{-1}$, память $20\,\text{с}$. (б)~$q=1$: $1$, $1$, $1\,\text{с}^{-1}$, $1\,\text{с}$. (в)~Память пропорциональна амплитуде шума: минута --- при шуме вшестеро сильнее. (г)~\nt{АЦЕТ} ниже в (а), выше в (б).}
\answer{ex:09:2}{(б)~$\hat s(t)=\frac1t\int_0^ty\,\dd t'$. (в)~$\frac12\sqrt{R/q}$ --- вдвое короче памяти. (г)~$t=\frac12\ln21\,\sqrt{R/q}\approx15\,\text{с}$.}
\answer{ex:09:3}{$\dd e/\dd t=-e/T+w-n/T$, дисперсия $qT/2+R/(2T)$; $T^*=\sqrt{R/q}$, дисперсия $\sqrt{qR}=P_\infty$; рост в $(\kappa+1/\kappa)/2$ раз: $1.25$ и $5.05$.}
\answer{ex:09:4}{(а)~$a_w=1-\theta/(mw)$: $0.22$ при $w=0.045$, $0.15$ при $w=0.041$. (б)~$a_c\approx0.38$ и $0.31$. (в)~$a<1-\theta/((2m-1)w)$: $0.59$ и $0.55$; границу чтения задаёт дописывание.}
\answer{ex:09:5}{Ошибки появляются при $M\approx40$--$60$, при $M=80$ --- $1.9$ лишних нейрона, при $M=100$ --- доля образа $0.86$; разброс $h$ совпадает с $0.92\sqrt{(M-1)p(1-p)/N}$ (при $M=80$: $0.172$ против $0.173$); $M_{\max}\approx80$, $M_{\max}\propto N/p$.}
\answer{ex:09:6}{Например, $\eta(a)=\eta\min\bigl(1,\max(0,(a-0.3)/0.6)\bigr)$. С $\eta a$ после $20$ чтений при $a=0.2$ все три чужих нейрона вошли в образ, при $a=0.3$ доля образа упала до $0.6$; с $\eta(a)$ --- ни одного, запись при $a\ge0.9$ та же (доля $1.00$).}
\answer{ex:09:7}{(а)~$40$ повторов; с $\eta(a)$ --- никогда. (б)~Ночью $a_B$ низкий, $a_S$ высокий: $S$ пишет, слушая $B$ как датчик. (в)~$200$ повторов, $10$ ночей; за день --- не больше $\approx M_{\max}/10$ образов.}
